\section{Bulk–boundary functoriality and spectral $R(z)$}
\label{sec:bulk-boundary-R}
\label{sec:spectral_braiding}
\label{sec:bulk-boundary}

\subsection{Boundary data and factorization module structure}
\label{subsec:boundary-module}
Let $A=(A_{\mathsf{ch}},A_{\mathsf{top}})$ be a $C_\ast\!\bigl(W(\mathsf{SC}^{\mathrm{ch,top}})\bigr)$–algebra.
A (topological) boundary condition along $t=0$ is modeled by a \emph{right $W(\mathsf{SC}^{\mathrm{ch,top}})$–module} $M$ supported on $\C\times \R_{\ge 0}$, i.e.\ a prefactorization algebra $\mathsf{Obs}^{\partial}$ on half‑rectangles with actions
\[
C_\ast\bigl(\FM_k(\C)\times E_1(m)\bigr)\otimes A_{\mathsf{ch}}^{\otimes k}\otimes M^{\otimes m}\longrightarrow M,
\]
compatible with $W$–coherence.  We refer to $M$ as a \emph{boundary factorization module}.

\subsection{Line operators and the spectral parameter}
\label{subsec:line-ops}
Boundary line operators $L$ are point‑supported defects along $t=0$ with insertion points $z\in\C$.
Given two parallel boundary lines $L_1(z_1)$ and $L_2(z_2)$ (with $z_1\neq z_2$), the bulk-to-boundary OPE in the strip determines a composition
\[
\mathsf{Comp}_{z_1,z_2} \colon L_1(z_1)\otimes L_2(z_2) \longrightarrow L_2(z_2)\otimes L_1(z_1),
\]
meromorphic in $z_{12}=z_1-z_2$ with logarithmic poles along $z_{12}=0$.
This composition is induced by the $W(\mathsf{SC}^{\mathrm{ch,top}})$–module structure and is computed by configuration integrals over $\FM_2(\C)$ with $E_1$–homotopies along the half‑interval.

\begin{definition}[Spectral braiding]
\label{def:spectral-braiding}
The \emph{spectral braiding} (or \emph{spectral $R$–matrix}) is the family of isomorphisms
\[
R_{L_1,L_2}(z)\colon L_1\otimes L_2 \xrightarrow{\ \sim\ } L_2\otimes L_1,
\qquad \text{for generic } z\in \C^\times \text{ (away from a discrete set of poles)},
\]
obtained from $\mathsf{Comp}_{z_1,z_2}$ by analytic continuation in $z=z_{12}$.
\end{definition}

\subsection{Yang–Baxter equation from $W$–coherence}
\label{subsec:YBE}
Consider three boundary lines $L_1,L_2,L_3$ at positions $z_1,z_2,z_3$.
The two ways to compose the spectral braidings $R_{12}(z_{12})$, $R_{13}(z_{13})$, $R_{23}(z_{23})$ correspond to different degenerations of $\FM_3(\C)$.

\begin{theorem}[Yang–Baxter equation]
\label{thm:YBE}
Assume {\rm (H1)--(H4)}.
The spectral braiding $R(z)$ satisfies the Yang–Baxter equation
\[
R_{12}(z_{12})\,R_{13}(z_{13})\,R_{23}(z_{23}) \;=\; R_{23}(z_{23})\,R_{13}(z_{13})\,R_{12}(z_{12})
\]
as morphisms $L_1\otimes L_2\otimes L_3 \to L_3\otimes L_2\otimes L_1$ for $z_i\neq z_j$.
\end{theorem}

\begin{proof}
We apply Stokes' theorem on $\FM_3(\C)$ to the logarithmic form $\omega_3^{\partial}$ governing the three-point bulk-to-boundary correlation of line operators $L_1(z_1), L_2(z_2), L_3(z_3)$.  Under (H3), $\omega_3^{\partial}$ extends to a logarithmic form on $\FM_3(\C)$ satisfying $d\omega_3^{\partial} = 0$ in the interior $\Conf_3(\C)$.

\medskip
\noindent\textbf{Step 1: Boundary strata of $\FM_3(\C)$.}
By the FM boundary structure (Definition~\ref{def:boundary_divisors}), the codimension-1 faces of $\partial\FM_3(\C)$ are the collision divisors
\[
D_{\{1,2\}},\quad D_{\{1,3\}},\quad D_{\{2,3\}},\quad D_{\{1,2,3\}}.
\]
The full-collision stratum $D_{\{1,2,3\}} \cong \FM_3(\C)$ contributes a term proportional to the identity (by (H4), the cluster limit where all three points collide simultaneously with no hierarchical ordering gives a subleading contribution that vanishes by degree counting: the weight form restricted to $D_{\{1,2,3\}}$ is exact on $\FM_3(\C)$).  Thus only the three pairwise divisors contribute.

\medskip
\noindent\textbf{Step 2: Spectral parameter assignments on each face.}
Each divisor $D_{\{i,j\}} \cong \FM_2(\C) \times \FM_2(\C)$ parametrizes configurations where $z_i$ and $z_j$ collide while the third point stays separated.  The cluster center inherits the spectral parameter of the pair, and the residue factorizes the three-point correlator into a product of two-point braidings.  Explicitly:

\begin{itemize}
\item \textbf{Face $D_{\{1,2\}}$}: Points $z_1, z_2$ collide first.  The inner braiding carries spectral parameter $z_{12} = z_1 - z_2$ and produces $R_{12}(z_{12})$.  The cluster center (at $z_2$) then braids with $z_3$, producing $R_{13}(z_{13})$ followed by $R_{23}(z_{23})$.  The total contribution is
\[
D_{\{1,2\}}: \quad R_{12}(z_{12})\, R_{13}(z_{13})\, R_{23}(z_{23}).
\]

\item \textbf{Face $D_{\{2,3\}}$}: Points $z_2, z_3$ collide first, giving inner braiding $R_{23}(z_{23})$.  The remaining braidings with $z_1$ yield $R_{13}(z_{13})$ then $R_{12}(z_{12})$.  The total contribution is
\[
D_{\{2,3\}}: \quad R_{23}(z_{23})\, R_{13}(z_{13})\, R_{12}(z_{12}).
\]

\item \textbf{Face $D_{\{1,3\}}$}: Points $z_1, z_3$ collide.  This face contributes
\[
D_{\{1,3\}}: \quad R_{13}(z_{13})\, R_{12}(z_{12})\, R_{23}(z_{23}).
\]
\end{itemize}

\noindent\textbf{Step 3: Stokes' theorem and corner cancellation.}
By Stokes' theorem on the manifold with corners $\FM_3(\C)$,
\[
0 = \int_{\FM_3(\C)} d\omega_3^{\partial} = \sum_{\{i,j\}} (-1)^{\sigma_{\{i,j\}}} \int_{D_{\{i,j\}}} \Res_{D_{\{i,j\}}}(\omega_3^{\partial}).
\]
Codimension-2 corners arise from nested collisions (e.g., $D_{\{1,2\}} \cap D_{\{2,3\}}$, where $z_1 \to z_2$ and $z_2 \to z_3$ simultaneously).  By Arnold cancellation (Lemma~\ref{lem:Arnold_cancellation}), the corner contributions cancel pairwise: for each pair of faces sharing a codimension-2 stratum, the two orderings of taking residues produce opposite signs via the Koszul sign rule on the normal covectors $d\varepsilon_{\{i,j\}}$.

\medskip
\noindent\textbf{Step 4: Assembling the identity.}
The orientation convention (Definition~\ref{def:orientation}) assigns signs $(-1)^{\sigma_{\{i,j\}}}$ to each face.  With the standard planar ordering $(1,2,3)$, the faces $D_{\{1,2\}}$ and $D_{\{2,3\}}$ carry opposite orientations.  The vanishing of the total boundary integral therefore gives
\[
R_{12}(z_{12})\,R_{13}(z_{13})\,R_{23}(z_{23}) \;-\; R_{23}(z_{23})\,R_{13}(z_{13})\,R_{12}(z_{12}) \;=\; 0,
\]
with the face $D_{\{1,3\}}$ not contributing an independent relation (its contribution is absorbed into the two terms above by the constraint $z_{13} = z_{12} + z_{23}$, which identifies the $D_{\{1,3\}}$ spectral parameter in terms of the other two).  This is the Yang--Baxter equation.
\end{proof}

\subsection{Classical limit and the infinitesimal $r(z)$}
\label{subsec:classical-r}
Let $\hbar$ denote the loop parameter of the perturbative HT theory.
Assume $R(z)$ admits an expansion
\[
R(z) \;=\; \mathrm{id} \,+\, \hbar\, r(z) \,+\, O(\hbar^2),
\]
in endomorphisms of $L_1\otimes L_2$.
Then $r(z)$ is a solution of the classical Yang–Baxter equation (CYBE) with spectral parameter:
\[
[ r_{12}(z_{12}), r_{13}(z_{13}) ] + [ r_{12}(z_{12}), r_{23}(z_{23}) ] + [ r_{13}(z_{13}), r_{23}(z_{23}) ] \;=\; 0.
\]

\begin{proposition}[Field–theoretic expression for $r(z)$]
\label{prop:field-theory-r}
Assume {\rm (H1)--(H4)}.
Let $\{\cdot{}_\lambda \cdot\}$ be the bulk $(-1)$–shifted $\lambda$–bracket on cohomology (Section~\ref{sec:Ainfty-to-PVA}). The kernel of the infinitesimal spectral $r$–matrix is computed at first order by
\[
r(z) \;\simeq\; \int_{0}^{\infty} \! \mathrm{d}\lambda\ \mathrm{e}^{-\lambda z}\ \bigl\{\; \cdot {}_\lambda \cdot \;\bigr\},
\]
under the BV pairing identifying bulk and boundary modes, and with the propagator fixing the time‑ordering on the half‑line.  The integral converges for $\operatorname{Re}(z) > 0$ and extends meromorphically to $\C^\times$. In the formal algebraic setting, this identity holds as an equality of formal power series in $z^{-1}$.
\end{proposition}

\begin{proof}
We prove both the Laplace-transform formula and the classical Yang--Baxter equation in four steps.

\medskip
\noindent\textbf{Step 1: Perturbative expansion.}
Write $R(z) = \id + \hh\, r(z) + O(\hh^2)$ in $\End(L_1 \otimes L_2)$.  At tree level ($\hh^0$), line operators compose freely and $R = \id$.  The first quantum correction $r(z)$ is computed by a single bulk-to-boundary propagator exchange: a bulk operator propagates from position $z_1$ on line $L_1$ to position $z_2$ on line $L_2$ (with $z = z_1 - z_2$) through the half-space $\C \times \R_{\geq 0}$.

\medskip
\noindent\textbf{Step 2: Propagator integral.}
Under (H2), the free propagator on $\C \times \R$ takes the form
\[
K(z,t) = \frac{\Theta(t)}{2\pi z}\, e^{-\mu(z) t} + (\text{regular}),
\]
where $\Theta(t)$ is the Heaviside function enforcing time-ordering and $\mu(z)$ encodes the holomorphic dependence.  The bulk $\lambda$-bracket kernel $K_\lambda$ is the Fourier-conjugate mode: $\{\cdot\,{}_\lambda\,\cdot\}$ acts on cohomology classes $[a], [b] \in H^\bullet(A,Q)$ via
\[
\{a_\lambda b\} = \int_0^\infty dt\; e^{\lambda t}\, K(z_{12}, t)\, \langle a \otimes b \rangle_{\mathrm{BV}},
\]
where $\langle \cdot \otimes \cdot \rangle_{\mathrm{BV}}$ is the BV pairing contracting bulk operators against boundary modes (cf.\ Section~\ref{sec:Ainfty-to-PVA}).

\medskip
\noindent\textbf{Step 3: Laplace transform.}
The first-order braiding $r(z)$ is obtained by evaluating the single-exchange diagram.  Performing the $t$-integration with the exponential weight $e^{-zt}$ from the spectral parameter gives
\[
r(z) = \int_0^\infty d\lambda\; e^{-\lambda z}\, \{\ \cdot{}_\lambda \cdot\ \}.
\]
This is the Laplace transform of the $\lambda$-bracket kernel, convergent for $\operatorname{Re}(z) > 0$ by the exponential decay of $K(z,t)$ in the topological direction (hypothesis (H2)).  Meromorphic continuation to $\C^\times$ follows from the meromorphy of the propagator in $z$.

For the free propagator $K(z,t) = \Theta(t)/(2\pi z)$, the $\lambda$-bracket is $\lambda$-independent and the Laplace integral evaluates to $r(z) \sim 1/z$ (consistent with the abelian CS result of Example~\ref{ex:Heisenberg_Yangian}).

\medskip
\noindent\textbf{Step 4: Classical Yang--Baxter equation and antisymmetry.}
Substituting $R(z) = \id + \hh\, r(z) + O(\hh^2)$ into the quantum YBE (Theorem~\ref{thm:YBE}), expanding to order $\hh^2$, and using $R_{ij} R_{ik} R_{jk} = R_{jk} R_{ik} R_{ij}$:
\begin{align*}
&(\id + \hh r_{12} + O(\hh^2))(\id + \hh r_{13} + O(\hh^2))(\id + \hh r_{23} + O(\hh^2)) \\
&\quad = (\id + \hh r_{23} + O(\hh^2))(\id + \hh r_{13} + O(\hh^2))(\id + \hh r_{12} + O(\hh^2)).
\end{align*}
At order $\hh^1$, both sides give $\id + \hh(r_{12} + r_{13} + r_{23})$, which is consistent.  At order $\hh^2$, the LHS contains $r_{12}r_{13} + r_{12}r_{23} + r_{13}r_{23}$ and the RHS contains $r_{23}r_{13} + r_{23}r_{12} + r_{13}r_{12}$.  Their difference yields
\[
[r_{12}(z_{12}),\, r_{13}(z_{13})] + [r_{12}(z_{12}),\, r_{23}(z_{23})] + [r_{13}(z_{13}),\, r_{23}(z_{23})] = 0,
\]
which is the classical Yang--Baxter equation with spectral parameter.

For antisymmetry, swapping $L_1 \leftrightarrow L_2$ sends $z \mapsto -z$ and exchanges tensor factors ($r \mapsto r_{21}$).  The spectral braiding satisfies $R_{12}(z) R_{21}(-z) = \id$ (inverse relation from the definition of braiding), so at order $\hh$:
\[
r_{12}(z) + r_{21}(-z) = 0, \qquad\text{i.e.,}\quad r(-z) = -r_{21}(z). \qedhere
\]
\end{proof}

\subsection{Koszul dual Hopf picture and factorisation quantum groups}
\label{subsec:Hopf}
Let $\mathbf B$ be the bar cooperad of $\mathsf{SC}^{\mathrm{ch,top}}$ and let $A$ be a $W(\mathsf{SC}^{\mathrm{ch,top}})$–algebra.
The bar construction $\mathsf{Bar}^{\mathrm{ch,top}}(A)$ is a colored coalgebra over $\mathbf B$; its cobar $\mathsf{Cobar}^{\mathrm{ch,top}}$ yields a (filtered) Hopf‑like object $\mathcal H$ encoding the OPE and higher products.

\begin{theorem}[Braiding on boundary categories]
\label{thm:braided-category}
Assume {\rm (H1)--(H4)}.
Let $\mathcal C_{\partial}$ be the (derived) category of boundary line operators with the tensor structure induced by the $W$–module structure of~$M$. Then $R(z)$ defines a meromorphic braiding on $\mathcal C_{\partial}$, and the classical limit $r(z)$ integrates to a filtered quasi‑triangular structure on the Koszul–dual Hopf object $\mathcal H$ (the "factorisation quantum group" in this setting).
\end{theorem}

\begin{proof}
We establish the meromorphic braiding and the quasi-triangular Hopf structure separately.

\medskip
\noindent\textbf{Part A: Meromorphic braiding on $\mathcal{C}_\partial$.}
The derived category $\mathcal{C}_\partial$ of boundary line operators inherits a monoidal structure from the $E_1$-composition along the boundary half-line $\R_{\geq 0}$: given right $W(\SCchtop)$-modules $M_1, M_2$ (boundary factorization modules in the sense of Section~\ref{subsec:boundary-module}), their tensor product $M_1 \otimes_{E_1} M_2$ is defined by the operadic composition in the open color.  The spectral braiding $R(z)$ of Definition~\ref{def:spectral-braiding} provides a natural isomorphism
\[
\beta_z \colon M_1 \otimes_{E_1} M_2 \xrightarrow{\ \sim\ } M_2 \otimes_{E_1} M_1
\]
for generic $z \in \C^\times$, arising from the exchange of insertion points in $\FM_2(\C)$.  By Theorem~\ref{thm:YBE}, $\beta_z$ satisfies the braid relation on triple tensor products, making $(\mathcal{C}_\partial, \otimes_{E_1}, \beta_z)$ a braided monoidal category with meromorphic dependence on $z$.

\medskip
\noindent\textbf{Part B: Filtered quasi-triangular structure on $\mathcal{H}$.}
Let $\mathcal{H} = \mathsf{Cobar}^{\mathrm{ch,top}}(\mathsf{Bar}^{\mathrm{ch,top}}(A))$ be the Hopf-like object from the bar--cobar construction.  By Theorem~\ref{thm:filtered-koszul}, $\mathcal{H}$ carries the holomorphic weight filtration $F^\bullet$ with $\gr^F \mathcal{H} \cong U(\mathfrak{g}[z])$ (the enveloping algebra of the loop algebra, on the closed color) tensored with an associative algebra (on the open color).

The classical $r$-matrix $r(z)$ defines an element $r(z) \in \mathcal{H} \widehat{\otimes}\, \mathcal{H}((z^{-1}))$ satisfying:
\begin{enumerate}[label=(\alph*)]
\item The CYBE (Proposition~\ref{prop:field-theory-r}, Step 4 of the proof), which is the infinitesimal form of the braid relation;
\item Antisymmetry $r(-z) = -r_{21}(z)$ (Theorem~\ref{thm:spectral_R_YBE}(iii));
\item The intertwining property $r(z)\,\Delta(x) - \Delta^{\mathrm{op}}(x)\,r(z) = 0$ at leading order in $\hh$ (from the quasi-triangularity established in the proof of Theorem~\ref{thm:Koszul_dual_Yangian}).
\end{enumerate}
These three properties define a \emph{filtered quasi-triangular structure}: $r(z)$ lives in filtration degree $F^1(\mathcal{H} \otimes \mathcal{H})$, and the CYBE holds modulo $F^2$.  The full quantum $R$-matrix $R(z) = \id + \hh\, r(z) + O(\hh^2)$ extends this to a quasi-triangular structure on the completed object $\widehat{\mathcal{H}}$, with the Yang--Baxter equation holding to all orders by Theorem~\ref{thm:YBE}.

The bar--cobar duality (Theorem~\ref{thm:bar-cobar-adjunction}) ensures that representations of $\mathcal{H}$ are equivalent to $\SCchtop$-modules, so the braided monoidal category $\mathcal{C}_\partial$ from Part A is identified with the category of $\mathcal{H}$-modules equipped with the braiding induced by $R(z)$.
\end{proof}

\subsection{OPE of Line Operators and the Spectral $R$-Matrix}

\subsubsection{Meromorphic Tensor Product}

When two line operators $\ell_1$ and $\ell_2$ are placed at positions $z_1, z_2 \in \C$ (separated in the holomorphic direction), their OPE is governed by a \textbf{meromorphic tensor product}:
\begin{equation}
\ell_1(z_1) \ell_2(z_2) \sim \sum_k C_k(z_1 - z_2) \, \ell_k,
\end{equation}
where $C_k(z)$ are meromorphic functions in $z$ encoding quantum corrections.

\begin{definition}[Spectral $R$-Matrix]
\label{def:spectral_R}
The \emph{spectral $R$-matrix} is a degree-1 element:
\begin{equation}
r(z) \in \mathcal{A}^! \widehat{\otimes} \mathcal{A}^!((z^{-1})),
\end{equation}
encoding the quantum corrections to the tensor product of line operators:
\begin{equation}
\ell_1(z) \otimes \ell_2(0) \sim e^{r(z)} (\ell_1 \otimes \ell_2).
\end{equation}
\end{definition}

\begin{remark}[Relation between $r(z)$ and $R(z)$]
The degree-1 element $r(z)$ is related to the braiding isomorphism $R(z)$ of Definition~\ref{def:spectral-braiding} by $R(z) = \exp(r(z))$ in the completed tensor product, where the exponential is well-defined by the nilpotency of the degree-1 part in the $A_\infty$ setting.
\end{remark}

\subsubsection{Properties of $r(z)$}

\begin{theorem}[Spectral $R$-Matrix Satisfies $A_\infty$ Yang--Baxter]
\label{thm:spectral_R_YBE}
Assume {\rm (H1)--(H4)}.
(Dimofte--Niu--Py \cite{DNP}, Theorem 5.2) The spectral $R$-matrix $r(z)$ satisfies:
\begin{enumerate}[label=(\roman*)]
\item \textbf{Maurer--Cartan equation}:
\begin{equation}
Q(r(z)) + \frac{1}{2}[r(z), r(z)] + \sum_{k \geq 3} m_k(r(z)^{\otimes k}) = 0;
\end{equation}

\item \textbf{$A_\infty$ Yang--Baxter equation}:
\begin{equation}
[r_{12}(w), r_{13}(z+w)] + [r_{12}(w), r_{23}(z)] + [r_{13}(z+w), r_{23}(z)] + (\text{higher $m_k$ terms}) = 0,
\end{equation}
in $\mathcal{A}^{\otimes 3}[[z^{-1}, w^{-1}, (z-w)^{-1}]]$;

\item \textbf{Skew-symmetry}:
\begin{equation}
r(-z) = -r_{21}(z),
\end{equation}
where $r_{21}$ means swapping the two tensor factors.
\end{enumerate}
\end{theorem}

\begin{proof}[Proof sketch]
\textbf{(i) Maurer--Cartan}: This follows from the associativity of the OPE. The OPE of three lines $\ell_1(z_1) \ell_2(z_2) \ell_3(z_3)$ must be independent of the order of expansion. This forces $r(z)$ to satisfy the MC equation in $\mathcal{A}^! \otimes \mathcal{A}^!$.

\textbf{(ii) $A_\infty$ Yang--Baxter}: The associativity of the OPE for three lines yields:
\begin{equation}
r_{23}(z) + (id \otimes \Delta_z)(r(z+w)) = r_{12}(w) + (\Delta_w \otimes id)(r(z)),
\end{equation}
where $\Delta_z : \mathcal{A}^! \to \mathcal{A}^! \otimes_{r(z)} \mathcal{A}^!$ is the coproduct deformed by $r(z)$. Expanding this identity and using the $A_\infty$ structure yields the $A_\infty$ Yang--Baxter equation.

\textbf{(iii) Skew-symmetry}: This follows from the symmetry of the tensor product: swapping $\ell_1 \leftrightarrow \ell_2$ corresponds to $z \to -z$ and exchanging tensor factors.
\end{proof}

\subsection{dg-Shifted Yangians}

\begin{definition}[dg-Shifted Yangian]
\label{def:dg_Yangian}
A \emph{dg-shifted Yangian} is an $A_\infty$ algebra $\mathcal{Y}$ equipped with:
\begin{enumerate}[label=(\roman*)]
\item A translation automorphism $\tau_z : \mathcal{Y} \to \mathcal{Y}[[z]]$, generated by a derivation $T$ of degree $(0, \text{odd}, 1)$:
\begin{equation}
\tau_z = \exp(z T) = \sum_{n \geq 0} \frac{z^n}{n!} T^n;
\end{equation}

\item A spectral $R$-matrix $r(z) \in \mathcal{Y} \widehat{\otimes} \mathcal{Y}[[z^{-1}, z]]$ of degree $(1, \text{odd}, 0)$;

\item A twisted coproduct $\Delta_z : \mathcal{Y} \to \mathcal{Y} \otimes_{r(z)} \mathcal{Y}[[z^{-1}, z]]$ (an $A_\infty$ algebra morphism);
\item A counit $\varepsilon : \mathcal{Y} \to \C$ (also an $A_\infty$ morphism);
\end{enumerate}
satisfying:
\begin{itemize}
\item Counit axioms: $\varepsilon(1) = 1$, $(\varepsilon \otimes id) \circ \Delta_z = id$, $(id \otimes \varepsilon) \circ \Delta_z \simeq \tau_z$;
\item Translation invariance: $\varepsilon(T(a)) = 0$, $\varepsilon(\tau_z(a)) = \varepsilon(a)$ for all $a \in \mathcal{Y}$;
\item Weak co-commutativity: $(\tau_w \otimes \tau_w) r(z) = r(z)$, $r(-z) = -r_{21}(z)$;
\item Co-associativity (quantum $A_\infty$ Yang--Baxter):
\begin{equation}
r_{23}(z) + (id \otimes \Delta_z)(r(z+w)) = r_{12}(w) + (\Delta_w \otimes id)(r(z)).
\end{equation}
\end{itemize}
\end{definition}
\begin{theorem}[Koszul Dual is a dg-Shifted Yangian]
\label{thm:Koszul_dual_Yangian}
(Dimofte--Niu--Py \cite{DNP}, Theorem 5.5) For a 3d HT gauge theory satisfying hypotheses H1--H4, the Koszul dual algebra $\mathcal{A}^!$ is a dg-shifted Yangian in the sense of Definition \ref{def:dg_Yangian}.
\end{theorem}

\begin{proof}
We verify each piece of structure in Definition~\ref{def:dg_Yangian}.

\medskip
\noindent\textbf{Step 1: Translation operator.}
The holomorphic translation $z \mapsto z + w$ acts on bulk local operators of the 3d HT theory.  Under the bar--cobar adjunction (Theorem~\ref{thm:bar-cobar-adjunction}), this induces a derivation $T$ of degree $(0,\text{odd},1)$ on the Koszul dual $\A^!$: if $\phi(z)$ is a boundary-localized field, then $T\phi = \partial_z \phi$.  The automorphism $\tau_w = \exp(wT) = \sum_{n \geq 0} w^n T^n / n!$ converges in $\A^![[w]]$ because $T$ raises filtration degree (the holomorphic weight filtration of Proposition~\ref{prop:gr-chiral}).

\medskip
\noindent\textbf{Step 2: Spectral $R$-matrix.}
The OPE of two parallel line operators $L_1(z_1)$, $L_2(z_2)$ is governed by configuration integrals over $\FM_2(\C) \times E_1(2)$.  Setting $z = z_1 - z_2$, the quantum corrections to the free tensor product define $r(z) \in \A^! \widehat{\otimes}\, \A^!((z^{-1}))$ as in Theorem~\ref{thm:spectral_R_YBE}.  The degree is $(1, \text{odd}, 0)$: degree 1 from the desuspension in the bar construction, odd under the $\Z/2$-grading from the fermion number, and weight 0 because the leading pole $1/z$ carries no holomorphic weight shift.

\medskip
\noindent\textbf{Step 3: Coproduct from the monoidal structure on $\mathcal{C}_{\mathrm{line}}$.}
Placing two parallel lines at $z_1, z_2 \in \C$ and taking the OPE defines the coproduct.  Concretely, for $x \in \A^!$, the element $\Delta_z(x)$ encodes how the action of $x$ on a single line decomposes when the line is resolved into a pair separated by spectral parameter $z$:
\[
\Delta_z \colon \A^! \longrightarrow \A^! \otimes_{r(z)} \A^![[z^{-1}, z]],
\]
where the twisted tensor product $\otimes_{r(z)}$ has differential $d_{\otimes} + [r(z), -]$.  This is an $A_\infty$ algebra morphism: the compatibility with higher operations $m_k$ follows from the $W(\SCchtop)$-module structure on boundary factorization modules (Section~\ref{subsec:boundary-module}), which ensures the OPE respects the full homotopy-coherent algebra structure.

The $R$-matrix intertwines the coproduct and its opposite:
\[
R(z)\,\Delta_z(x) = \Delta_z^{\mathrm{op}}(x)\, R(z) \qquad \text{for all } x \in \A^!.
\]
This is the defining property of a quasi-triangular structure.  It follows from the definition of $R(z)$ as the braiding isomorphism (Definition~\ref{def:spectral-braiding}): applying $R(z)$ swaps the two tensor factors, converting $\Delta_z$ (first factor at $z_1$, second at $z_2$) to $\Delta_z^{\mathrm{op}}$ (factors swapped).

\medskip
\noindent\textbf{Step 4: Counit.}
Define $\varepsilon \colon \A^! \to \C$ by the augmentation: removing a line from a configuration corresponds to the counit of the $E_1$-coalgebra structure.  Explicitly, $\varepsilon$ evaluates a line operator on the vacuum module.  Then:
\begin{itemize}
\item $\varepsilon(1) = 1$ (the identity line evaluates to the scalar identity);
\item $(\varepsilon \otimes \id) \circ \Delta_z = \id$ (removing the first line recovers the second);
\item $(\id \otimes \varepsilon) \circ \Delta_z \simeq \tau_z$ (removing the second line at position $z_2$ and keeping the first at $z_1$ shifts by $z = z_1 - z_2$ via holomorphic translation).
\end{itemize}

\medskip
\noindent\textbf{Step 5: Translation invariance.}
Simultaneous holomorphic translation $z_i \mapsto z_i + w$ for all lines preserves the spectral differences $z_{ij}$.  Therefore:
\begin{itemize}
\item $\varepsilon(T(a)) = \varepsilon(\partial_z a) = 0$ (the augmentation kills derivatives, since the vacuum is translation-invariant);
\item $(\tau_w \otimes \tau_w)\, r(z) = r(z)$ (the $R$-matrix depends only on $z = z_1 - z_2$, not on the individual positions).
\end{itemize}

\medskip
\noindent\textbf{Step 6: Co-associativity (quantum $A_\infty$ Yang--Baxter).}
The identity
\[
r_{23}(z) + (\id \otimes \Delta_z)(r(z+w)) = r_{12}(w) + (\Delta_w \otimes \id)(r(z))
\]
expresses the consistency of the three-line OPE: resolving the pair $(L_2, L_3)$ first and then braiding $L_1$ past the pair (LHS) must agree with braiding $L_1$ past $L_2$ first and then resolving (RHS).  This is a rewriting of the Yang--Baxter equation (Theorem~\ref{thm:YBE}) in coproduct language, obtained by substituting $R(z) = \exp(r(z))$ and expanding the three-fold braiding identity to leading nontrivial order.

\medskip
\noindent\textbf{Step 7: Yangian structure.}
The spectral parameter $z$ makes this a \emph{Yangian} rather than a mere quasi-triangular Hopf algebra: $\Delta_z$, $r(z)$, and $\tau_z$ assemble into a deformation of $U(\mathfrak{g}[z])$ (the universal enveloping algebra of the loop algebra).  In the associated graded with respect to the holomorphic filtration, $r(z) \sim r_0/z$ (simple pole) and $\Delta_z$ reduces to the standard Yangian coproduct $\Delta_z(x) = x \otimes 1 + 1 \otimes x + O(z^{-1})$.  This matches Definition~\ref{def:dg_Yangian}.
\end{proof}

\subsection{Connection to Latyntsev's Factorisation Quantum Groups}

The spectral $R$-matrix $r(z)$ and the meromorphic tensor product of line operators fit naturally into the framework of \textbf{factorisation quantum groups} developed by Latyntsev \cite{Lat}.

\begin{definition}[Factorisation Quantum Group (Informal)]
\label{def:factorisation_QG}
A \emph{factorisation quantum group} is a categorical structure encoding:
\begin{itemize}
\item A monoidal category $\mathcal{C}$ of representations;
\item A meromorphic tensor product $\otimes_z : \mathcal{C} \times \mathcal{C} \to \mathcal{C}[[z^{-1}, z]]$, parametrized by a spectral parameter $z \in \C$;
\item An $R$-matrix $R(z) : V \otimes_z W \to W \otimes_{-z} V$ satisfying the quantum Yang--Baxter equation.
\end{itemize}
\end{definition}

\begin{theorem}[Line Operators Form a Factorisation Quantum Group]
\label{thm:lines_factorisationQG}
Assume {\rm (H1)--(H4)}.
The category $\mathcal{C}_{\text{line}}$ of line operators in a 3d HT theory, equipped with the meromorphic OPE and spectral $R$-matrix $r(z)$, forms a factorisation quantum group in the sense of Latyntsev \cite{Lat}.
\end{theorem}

\begin{proof}
We verify the three axioms of a factorisation quantum group (Definition~\ref{def:factorisation_QG}).

\medskip
\noindent\textbf{(i) Monoidal structure.}
The category $\mathcal{C}_{\mathrm{line}}$ has objects given by boundary line operators and morphisms given by boundary-localized local operators (intertwiners).  The monoidal product is concatenation in the $E_1$-direction along the boundary $t = 0$: given lines $L_1, L_2$ placed at distinct points in $\C$, the product $L_1 \otimes L_2$ is defined by the $W(\SCchtop)$-module structure (Section~\ref{subsec:boundary-module}).  Associativity and unit constraints follow from the $E_1$-coherence encoded in the $W$-resolution.  By Theorem~\ref{thm:lines_as_modules}, $\mathcal{C}_{\mathrm{line}}$ is equivalent (as a dg category) to the category of $\A^!$-modules.

\medskip
\noindent\textbf{(ii) Meromorphic tensor product.}
For each $z \in \C^\times$, the OPE of line operators at separation $z$ defines a meromorphic family of bifunctors
\[
\otimes_z \colon \mathcal{C}_{\mathrm{line}} \times \mathcal{C}_{\mathrm{line}} \longrightarrow \mathcal{C}_{\mathrm{line}}[[z^{-1}, z]].
\]
Concretely, $L_1 \otimes_z L_2$ is the line operator obtained by placing $L_1$ at position $z$ and $L_2$ at position $0$ and taking the bulk-to-boundary OPE.  The OPE coefficients $C_k(z)$ are meromorphic in $z$ with poles only at $z = 0$ (collision singularities), as guaranteed by (H2).  The family $\otimes_z$ varies holomorphically in $z$ on $\C^\times$: by (H3), the FM-compactified configuration integrals that define the OPE are smooth functions of the spectral parameter away from the collision divisor.

\medskip
\noindent\textbf{(iii) Braiding by $R(z)$.}
The spectral braiding $R(z) \colon L_1 \otimes_z L_2 \xrightarrow{\sim} L_2 \otimes_{-z} L_1$ (Definition~\ref{def:spectral-braiding}) provides a natural isomorphism between $\otimes_z$ and the opposite product $\otimes_{-z}^{\mathrm{op}}$.  By Theorem~\ref{thm:YBE}, $R(z)$ satisfies the quantum Yang--Baxter equation, which is precisely the braid relation required for a braided monoidal category with spectral parameter.

\medskip
\noindent\textbf{(iv) Factorisation property.}
The distinguishing feature of a \emph{factorisation} quantum group (as opposed to an ordinary quantum group) is that the braiding depends holomorphically on $z$, upgrading the braided monoidal category to a \emph{factorisation category} on $\C$.  In Latyntsev's framework \cite{Lat}, this means: the assignment $z \mapsto (\mathcal{C}_{\mathrm{line}}, \otimes_z, R(z))$ defines a constructible cosheaf on the Ran space $\Ran(\C)$.  For our category, this factorisation structure is inherited from the prefactorization algebra $\Obs^{\partial}$ on $\C \times \R_{\geq 0}$: restricting $\Obs^{\partial}$ to disks at different points of $\C$ and taking the $E_1$-coinvariants in the $\R$-direction yields fibers of the cosheaf, and the factorisation isomorphisms come from the prefactorization algebra structure maps for disjoint disks.

The holomorphic variation of the braiding in $z$ is therefore a direct consequence of the holomorphic-topological factorisation algebra structure of the 3d HT theory, restricted to the boundary.
\end{proof}

\subsection{Example: Abelian CS and Heisenberg Yangian}

\begin{example}[Heisenberg Yangian from $U(1)$ CS]
\label{ex:Heisenberg_Yangian}
For abelian $U(1)$ Chern--Simons theory at level $k$ (Example \ref{ex:abelian_CS_complete}):
\begin{itemize}
\item Bulk algebra: $\mathcal{A}_{\text{bulk}} = \widehat{\mathfrak{u}(1)}_k$ (affine Kac--Moody);
\item Koszul dual: $\mathcal{A}^! = Y(\mathfrak{u}(1))$ (Heisenberg Yangian).
\end{itemize}

The spectral $R$-matrix is:
\begin{equation}
r(z) = k \, \frac{J \otimes J}{z} + \text{(higher terms)}.
\end{equation}

This satisfies the classical Yang--Baxter equation:
\begin{equation}
[r_{12}(w), r_{13}(z+w)] + [r_{12}(w), r_{23}(z)] + [r_{13}(z+w), r_{23}(z)] = 0.
\end{equation}

For non-abelian gauge groups $G$, the Koszul dual becomes the \textbf{Yangian $Y(\mathfrak{g})$} or \textbf{quantum affine algebra $U_q(\widehat{\mathfrak{g}})$}, depending on the CS level and boundary conditions.
\end{example}

%%%%%%%%%%%%%%%%%%%%%%%%%%%%%%%%%%%%%%%%%%%%%%%%%%%%%%%%%%%%%%%%%%%%%%%%%%%%%
% SECTION 9 — WORKED EXAMPLES
%%%%%%%%%%%%%%%%%%%%%%%%%%%%%%%%%%%%%%%%%%%%%%%%%%%%%%%%%%%%%%%%%%%%%%%%%%%%%
